\section{Conclusion}

The central finding of this paper is that host-level hygiene posture, patch debt, AD privilege exposure, and telemetry freshness, adds substantial prioritization value over EPSS-only rankings when operationalized as a Hygiene Risk Score and integrated into a pre-calibrated weighted scorer. On a pre-registered, fully reproducible benchmark over the EEHDA fleet, HygienePrio-full outperforms EPSS-only by 31 pp at P@50 across all 25 of 25 evaluation seeds (non-overlapping BCa 95\% CIs). Under real EPSS/KEV attribute distributions the aggregate gap is +2.4 pp (\S\ref{sec:external_validity}), while host-side discrimination at major-CVE disclosure moments, exactly when remediation capacity is scarcest, remains the operationally significant advantage (\S\ref{sec:retrospective}). The practical take-home message is concrete: of the three hygiene dimensions, patch posture alone accounts for $\approx$20 pp of this improvement. Operations teams with limited instrumentation capacity should prioritize patch posture tracking above more complex hygiene signals such as AD group exposure or telemetry freshness monitoring.

This result completes a three-paper methodological arc. The VulnPrio study (VulnPrio) established that CVE-level feature augmentation does not improve over EPSS-only, explicitly attributing the null to the absence of dynamic host-level signals. The HygieneBench benchmark (HygieneBench) demonstrated that patch posture and AD group drift carry discriminative structure in fleet telemetry. HygienePrio operationalizes these signals in a remediation prioritization scorer. The progression from honest null through diagnostic signal identification to positive integration result illustrates how pre-registered negative results generate scientific value by constraining and directing subsequent work.

The headline magnitudes are benchmark measurements. The ground-truth definition incorporates HRS as a labelling condition, which contributes a structural component to the benchmark gap that real-exploitation-data evaluation would remove; the external-validity and retrospective analyses (\S\S\ref{sec:external_validity}, \ref{sec:retrospective}) quantify and bound this component. The calibrated weights were estimated on five seeds and should be re-fitted before real-fleet deployment. We therefore present HygienePrio as a validated benchmark contribution with a characterized boundary of applicability, ahead of prospective real-fleet evaluation.

\textbf{Reproducibility.} The evaluation code, fleet generator, frozen results, and pre-registration protocol are available in the accompanying repository at \url{https://github.com/Harshavardhanmalla-Labs/research-papers/tree/main/paper4}. The real-data artifacts referenced in \S\S\ref{sec:external_validity} to \ref{sec:retrospective} live under \texttt{real\_data/} at the repository root (not under \texttt{paper4/}).

\textbf{Future work.} The most important next step is evaluation on real, appropriately anonymized fleet telemetry using exploitation event logs (not HRS) as ground truth. Secondary directions include: (i) extending HRS with network-layer signals (firewall coverage, internet-facing service flags) to test whether additional dimensions contribute beyond the current three; (ii) evaluating HygienePrio under rolling daily EPSS score updates rather than static per-seed snapshots; and (iii) testing whether non-linear scoring functions trained on historical remediation outcomes can outperform the pre-calibrated linear scorer, which is the binding simplification assumption of this work.
